\documentclass[10pt,a4paper]{amsart}
\usepackage[margin=19mm]{geometry}
\usepackage{amsmath,amssymb,mathtools}
\usepackage[scr=rsfs,cal=euler]{mathalfa}
\usepackage{xfrac}
\usepackage[unicode,hidelinks,bookmarks=false]{hyperref}
\newcommand{\ie}{i.e.\ }
\newcommand{\cf}{cf.\ }
\newcommand{\resp}{respectively\ }
\newcommand{\Subsection}{\subsection}
\providecommand{\nobreakdash}{\nobreak\hbox{-}}
\setlength{\emergencystretch}{2em}
\multlinegap0pt
\usepackage{bm}
\usepackage{bbm}
\usepackage{dsfont}
\usepackage{xspace}
\usepackage{cancel}
\usepackage{mathtools}
\usepackage{amsthm}
\makeatletter
\@ifundefined{theorem}{\newtheorem{theorem}{Theorem}}{}
\makeatother
\title[Continuity of Functions on Bare Representation of Graphs und]{Continuity of Functions on Bare Representation of Graphs under Star Topology}
\author{Rodolfo E. Maza}
\date{Source: arXiv:2507.13968v1 (CC BY 4.0)}
\begin{document}
\maketitle

\noindent\emph{Source and licence.} Continuity of Functions on Bare Representation of Graphs under Star Topology. Version arXiv:2507.13968v1 (CC BY 4.0), distributed under the Creative Commons Attribution 4.0 International licence, as recorded in the arXiv licence metadata for this exact version and shown on the abstract page https://arxiv.org/abs/2507.13968v1. Full rights notice: licenses/arxiv-2507.13968-CC-BY-4.0.txt.\par\smallskip

\noindent\textit{Editorial scope.} Counted unit: the Theorem whose source label is \texttt{None} in the article (source0.tex lines 367--369, source label \texttt{None}), delivered together with the article's own complete original proof of it (source0.tex lines 370--411, one \texttt{proof} environment of 42 lines). No mathematics is written, repaired or re-derived: statement and proof are the article's own text line for line. Required context retained uncounted: none. Every definition, display and label of those lines is retained unchanged. Omitted from this package and not redistributed: 1--366, 412--929. Nothing inside a retained excerpt is omitted or altered: every retained body is delivered exactly as the article's lines are written, so no omission receipt of any kind accompanies this package. Citations: the retained excerpts carry no citation command at all, so no bibliography entry of the article is transcribed and no reference is redistributed. Reference adaptation: none was needed and none was made; every reference inside the delivered unit resolves inside it, so no unresolved reference or citation is produced. Printed slips kept literal: no printed word, symbol, hypothesis or number of the counted statement or of its proof was altered. Layout and class: the article's own class is \texttt{amsart} with packages including graphicx, tikz, tikz-cd; the wrapper is the lane's pinned \texttt{amsart} editorial wrapper at 10pt on A4, which restates the theorem-like environments the retained text needs and adds the article's own macro definitions that the retained text needs (none), with extra wrapper packages (bm, bbm, dsfont, xspace, cancel, mathtools). The delivered numbering is the unit's own. Assets: no retained span contains a figure environment, a graphics-inclusion command, a PDF-page inclusion, a table or a TikZ picture; no image file is copied into this package, the package declares no asset dependency and no image file is redistributed with it. Licence: version arXiv:2507.13968v1 of this article is distributed under the Creative Commons Attribution 4.0 International licence, as recorded in the arXiv licence metadata for this exact version and shown on the abstract page https://arxiv.org/abs/2507.13968v1. A full rights notice is delivered at licenses/arxiv-2507.13968-CC-BY-4.0.txt. Characters: every retained excerpt is pure ASCII, so the delivered engine, XeLaTeX with its default Latin Modern text font, reports no missing character.
	\begin{theorem}
		A graph \(G\) is connected if and only if \(B(G)\) is connected.
	\end{theorem}

	\begin{proof}
		Assume that \( G \) is a connected graph. We aim to show that its Bare representation
		\(
		B(G) \) under
		the star topology is also connected.
		
		For contradiction, suppose that \( B(G) \) is disconnected. Then, there exist two
		non-empty disjoint
		open sets \( U \) and \( V \) in \( B(G) \) such that:
		\[
		B(G) = U \cup V
		\]
		Choose a vertex \( u \in U \) and another vertex \( v \in V \). Since \( G \) is
		connected, there
		exists a finite path \( P: u = u_1u_2 \cdots u_n = v \).
		
		Consider each edge \( e = u_iu_{i+1} \) in the path. If one endpoint of an edge
		is
		in \( U \), then
		by the definition of open sets in the star topology, the entire neighborhood star
		around that vertex
		(including the edge itself) must lie within \( U \). Similarly, if a vertex is in
		\(
		V \), its
		neighborhood star lies entirely in \( V \).
		
		However, along the path from \( u \) to \( v \), there must be an edge where one
		endpoint
		transitions from \( U \) to \( V \). This transition would imply that the edge
		belongs to both \( U
		\) and \( V \), contradicting their disjointness. Hence, \( B(G) \) is connected.
		
		To establish the converse, we assume that \( B(G) \) is connected under the star
		topology and demonstrate that \( G \) is connected. Assume, for contradiction, that \( G \) is not connected. Then, it can be partitioned into two vertex-disjoint non-empty
		subgraphs \( H_1 \) and \( H_2 \). Also, \( H_1 \) and \( H_2 \) are edge-disjoint.
		
		Since the set \( B(H_i) \) consists of all vertices and edges of \( H_i \), it follows that \( B(H_1) \) and \( B(H_2) \) are non-empty disjoint sets whose union is $B(G)$. In the star topology, any neighborhood star around a vertex	within \( H_1 \) includes only vertices and edges from \( H_1 \), meaning \( B(H_1) \) is an open set. The same applies to \( B(H_2) \). Therefore, we have:
		\[
		B(G) = B(H_1) \cup B(H_2)
		\]
		where both \( B(H_1) \) and \( B(H_2) \) are open, non-empty, and disjoint. This contradicts the assumption that \( B(G) \) is connected. Hence, our initial assumption that \( G \) is disconnected must be false, which shows that \( G \) is connected.
	\end{proof}

\end{document}
